% ============================================================
\section{Conclusion}
% ============================================================
\label{sec:conclusion}

This work addresses an open problem in LLM deployment security: while \citet{staab2024beyond} established that modern LLMs can infer sensitive personal attributes from seemingly innocuous user-generated text, the question of how to \emph{detect and block} such inference requests at deployment time had not been studied.

Over a benchmark of 10,800 adversarial records (4 open-weight LLMs $\times$ 100 SynthPAI profiles $\times$ 9 UAL attack variants $\times$ 3 defense conditions) and 9,600 benign records (4 models $\times$ 100 profiles $\times$ 8 tasks $\times$ 3 defense conditions), our evaluation of Prompt Guard~2---the current state of practice for input-side LLM defense---demonstrates a systematic and structural failure against UAL-Inference requests expressed in natural language. Only two of nine attack variants are correctly detected, yielding an overall detection rate of 46.2\%. The most evasive formulation (\texttt{ual\_inference\_evasive\_stealth}) receives a 99.4\% \textsc{Safe} confidence score per query despite encoding a complete profiling intent. Furthermore, Prompt Guard~2 introduces a 50\% false-positive rate on benign tasks, entirely blocking four of eight standard utility task categories (Summary, Reply Draft, Translation, Keyword Extraction). The root cause is a structural misalignment: Prompt Guard~2 evaluates the \emph{syntax} of a prompt---injection-style markers, structured output constraints---rather than its \emph{semantic purpose}.

We operationalize the necessary shift through \textbf{UAL Semantic Guard}, which directly answers the open problem: it detects whether a prompt's underlying purpose is User Attribute Inference, regardless of its phrasing. Implemented as a zero-shot LLM-as-a-Judge, UAL Semantic Guard correctly identifies all nine UAL formulations across all four models, all eight personal attribute categories, and all five SynthPAI hardness levels---achieving 100\% detection---while producing 38 false positives across 3,200 benign evaluations (1.19\% FPR), concentrated primarily on a single benign task (Sentiment Analysis). This places UAL Semantic Guard in a near-ideal position relative to Prompt Guard~2 in the security-utility tradeoff space.

The main limitation of the current LLM-as-a-Judge implementation is its operational cost on benign queries: a mean pre-classification overhead of $+2.0$\,s ($+42.1\%$) and $+16.6\%$ energy per benign query compared to unguarded execution. This overhead makes the current implementation most appropriate for asynchronous or low-throughput pipelines. A lightweight fine-tuned classifier variant is the primary direction for closing the latency-security gap.

More broadly, our findings underscore a fundamental tension in the deployment of capable LLMs: the semantic inference abilities that make these systems valuable also render them effective profiling instruments when redirected by adversarial prompts. There is no version of this capability that can be selectively enabled for legitimate analytical use and simultaneously rendered inert for attribute inference---the underlying computation is the same. This implies that defenses must always operate at the level of \emph{intent}, not \emph{capability}.

Future work includes: scaling the evaluation to the full SynthPAI corpus and multilingual settings; training a lightweight fine-tuned classifier variant of UAL Semantic Guard to reduce the latency overhead while preserving semantic detection; extending the threat model to cover multi-turn splitting attacks; and evaluating robustness against adaptive adversaries who iteratively refine evasion strategies upon observing Judge verdicts.

We will release all evaluation scripts, prompt templates, and the UAL Semantic Guard implementation as open-source artifacts to support reproducibility and continued research on LLM privacy defense.
